% ============================================================

%  2. LITERATURE REVIEW

% ============================================================

\chapter{Literature Review}\label{sec:literature_review}

A natural starting point is the Phillips curve, which relates inflationary pressure to labour-market slack. In its simplest form, lower unemployment is expected to increase wage and price pressure, thereby raising inflation. However, the empirical usefulness of the Phillips curve for forecasting has been widely debated. While the relationship remains theoretically important, its forecasting performance has been shown to vary across countries, time periods, and inflation regimes \parencite{stock_watson_2008, stock2019}.

A key challenge in inflation forecasting is that inflation often displays strong serial dependence. Following the terminology in \textit{Forecasting: Principles and Practice}, ARIMA models provide a natural benchmark because they are designed to describe the autocorrelations in a time series. In this framework, differencing is used only when needed to obtain stationarity, while autoregressive and moving average terms capture the remaining dependence in the stationary series. The original series is therefore not required to be white noise. Instead, a well specified forecasting model should leave innovation residuals that are approximately uncorrelated and centred around zero.


This makes univariate ARIMA models useful benchmark models for inflation forecasting. They do not impose a structural economic relationship, but instead use the information contained in the past values of inflation itself. Such a benchmark is important because any economically motivated model should improve upon the forecasts generated by the inflation series' own history. If labour market variables do not improve forecast accuracy relative to this benchmark, this suggests that the Phillips curve has limited practical forecasting value in the given sample.

Dynamic regression models, or regression models with ARIMA errors, provide a way to combine economic predictors with time-series dynamics. In a standard regression model, the error term is assumed to be white noise. In time-series applications, this assumption is often unrealistic because the regression errors may still be autocorrelated. Regression with ARIMA errors addresses this by allowing inflation to depend on explanatory variables, such as unemployment, while modelling the remaining serial correlation in the error term using an ARIMA process. This approach is therefore well suited to testing whether unemployment contains predictive information for inflation beyond what is already captured by the autocorrelation structure of inflation itself.

{\color{red}

Vector autoregressions (VARs) provide an alternative multivariate framework when several macroeconomic variables are treated as jointly endogenous. A VAR model allows Danish inflation, Danish unemployment, and European inflation to interact through their own lags and each other's lags. This can be useful when the purpose is to analyse feedback effects and cross-variable dynamics rather than modelling Danish inflation conditional on selected predictors. However, VAR models can become heavily parameterised in relatively small samples, making lag selection and model reduction important practical issues.
}

For a small open economy such as Denmark, broader European inflation dynamics may also be relevant. Danish prices can be affected by import prices, energy markets, exchange-rate arrangements, and monetary conditions in the euro area. Including European macroeconomic information may therefore improve inflation forecasts, but it also increases the risk of overfitting. General-to-Specific (GETS) modelling offers a systematic way to reduce a larger dynamic specification by retaining the most statistically relevant variables and lags \parencite{doornik2014}.

Against this background, the present paper compares a univariate ARIMA benchmark with a dynamic regression model inspired by the Phillips curve. The benchmark model captures the persistence in Danish inflation using only its own history, whereas the dynamic regression model incorporates labour-market information and allows the remaining errors to follow an ARIMA process. This makes it possible to assess whether unemployment contains predictive information for Danish inflation beyond what is already captured by the inflation series itself.

{\color{red}

If a VAR specification is also included, the analysis further examines whether joint dynamics between Danish inflation, Danish unemployment, and European inflation improve forecast performance relative to the univariate benchmark.

}